\chapter{Реализация и тестирование}
\label{chap:implementation-and-testing}

\section{Структура библиотеки}

Реализация библиотеки сосредоточена в каталоге `crates/my-infra-otel/src`, где
размещены модули конфигурации, инициализации, серверной и клиентской
трассировки, журналирования, записи прикладных событий и обработки ошибок.
Отдельный каталог `crates/my-infra-otel/tests` содержит проверки публичного
интерфейса и наблюдаемого поведения. Практическое подключение библиотеки
демонстрируется в каталоге `examples`, а инфраструктурный контур для экспорта
трейсов в Jaeger описан в каталоге `infra`.

\begin{verbatim}
crates/my-infra-otel/src/
  client.rs
  config.rs
  error.rs
  events.rs
  guard.rs
  header_capture.rs
  init.rs
  labels.rs
  layer.rs
  lib.rs
  logging.rs
  propagation.rs
  request_builder.rs
  internal/
    mod.rs
    otel.rs
    subscriber.rs
    time.rs

crates/my-infra-otel/tests/
  config.rs
  header_attrs.rs
  header_capture.rs
  log_correlation.rs
  propagation.rs
  shutdown.rs

examples/
  checkout-api/
  order-processor/
  risk-engine/
  demo-ui/

infra/
  docker-compose.yml
  otel-collector-config.yml
\end{verbatim}

Корневой модуль `lib.rs` формирует публичный контракт библиотеки и переэкспортирует
основные типы: `TracingConfig`, `TracingGuard`, `MyOtelTracingLayer`,
`TracedHttpClient`, `record_event` и связанные вспомогательные сущности. Такое
разделение позволяет отделить стабильный интерфейс для прикладного кода от
внутренней инфраструктурной реализации.

\section{Реализация промежуточного слоя}

Серверная обработка входящих HTTP-запросов реализована в модуле `layer.rs`, где
объявлены `MyOtelTracingLayer`, его построитель и внутренняя обертка
`MyOtelTracingService<S>`. Модуль использует стандартные абстракции Tower
`Layer` и `Service`, что обеспечивает совместимость с Axum и позволяет
подключать библиотеку на уровне маршрутизатора~\cite{towerDocs,towerServiceDocs,towerLayerDocs,axumDocs}.

При обработке входящего запроса промежуточный слой:
\begin{enumerate}
    \item определяет метод, путь и маршрут запроса;
    \item создает серверный спан `http.request`;
    \item при наличии маршрута записывает его в `http.route`;
    \item извлекает родительский контекст из входящих заголовков;
    \item при необходимости захватывает разрешенные заголовки и агент клиента;
    \item запускает внутренний обработчик внутри созданного спана;
    \item после завершения обработки фиксирует код ответа, длительность и тип
    ошибки.
\end{enumerate}

Характерный фрагмент этой реализации приведен ниже:

\begin{verbatim}
let span = tracing::info_span!(
    "http.request",
    otel.kind = "server",
    http.request.method = %method,
    url.path = %path,
    http.route = tracing::field::Empty,
    http.response.status_code = tracing::field::Empty,
    duration.ms = tracing::field::Empty,
    error = tracing::field::Empty,
    error.type = tracing::field::Empty,
    otel.status_code = tracing::field::Empty,
    otel.status_message = tracing::field::Empty,
);
\end{verbatim}

Далее обработчик вызывается через `instrument(span.clone())`, благодаря чему
вложенные спаны и прикладные события автоматически попадают в тот же контекст
запроса. Если ответ имеет код ошибки, библиотека записывает признаки ошибки как
в полях спана `tracing`, так и в атрибутах OpenTelemetry.

\section{Реализация передачи контекста}

Передача контекста трассировки реализована в модуле `propagation.rs`. Этот
модуль не вводит собственного протокола, а адаптирует `http::HeaderMap` к
стандартным интерфейсам OpenTelemetry для извлечения и внедрения данных
трассировки~\cite{w3cTraceContext,opentelemetryPropagationDocs}.

В модуле определены две основные функции:

\begin{verbatim}
pub(crate) fn inject_context(
    context: &opentelemetry::Context,
    headers: &mut HeaderMap
)

pub(crate) fn extract_context(
    headers: &HeaderMap
) -> opentelemetry::Context
\end{verbatim}

Для их работы используются два адаптера:
\begin{enumerate}
    \item `HeaderExtractor`, читающий значения заголовков;
    \item `HeaderInjector`, записывающий значения заголовков.
\end{enumerate}

На практике это означает, что входящий запрос может продолжить уже существующую
трейс, если заголовок `traceparent` пришел от внешнего клиента или другого
сервиса, а исходящий запрос автоматически передает свой контекст следующему
участнику цепочки. Таким образом, библиотека обеспечивает совместимость с W3C
Trace Context на обеих границах HTTP-взаимодействия.

\section{Реализация экспортера OTLP}

Запуск подсистемы трассировки сосредоточен в модулях `init.rs`,
`internal/otel.rs` и `guard.rs`. Публичная функция `init_global_tracing`
принимает `TracingConfig`, предотвращает повторную инициализацию через атомарный
флаг и затем строит общий конвейер трассировки.

Ключевую инфраструктурную работу выполняет модуль `internal::otel`. Он:
\begin{enumerate}
    \item регистрирует распространитель W3C Trace Context;
    \item собирает ресурсные атрибуты сервиса, включая `service.name`,
    `service.version` и `deployment.environment`;
    \item настраивает экспортер OTLP/HTTP по адресу из конфигурации;
    \item выбирает пакетный или простой режим установки в зависимости от
    наличия активного `tokio`-исполнителя.
\end{enumerate}

Существенный фрагмент настройки экспортера имеет следующий вид:

\begin{verbatim}
let pipeline = opentelemetry_otlp::new_pipeline()
    .tracing()
    .with_exporter(
        opentelemetry_otlp::new_exporter()
            .http()
            .with_endpoint(config.otlp_endpoint.clone())
            .with_timeout(config.export_timeout),
    );
\end{verbatim}

Управляемое завершение работы реализовано типом `TracingGuard`. При вызове
`shutdown()` он выполняет завершение поставщика трассировщиков в отдельном потоке
и ограничивает ожидание тайм-аутом из конфигурации. Такое решение снижает риск
потери уже накопленных спанов при остановке процесса и не позволяет приложению
зависнуть на неопределенное время при недоступности внешнего приемника.

Для практической передачи трейсов библиотека использует внешний OpenTelemetry
Collector. В файле `infra/otel-collector-config.yml` описан приемник OTLP по
HTTP на порту `4318`, пакетная обработка и дальнейшая отправка данных в
Jaeger~\cite{otlpSpec,otlpExporterSpec,jaegerDocs}. В файле
`infra/docker-compose.yml` всем демонстрационным сервисам передается переменная
`OTLP_ENDPOINT=http://otel-collector:4318/v1/traces`, что обеспечивает единый
маршрут экспорта.

\section{Структурированные журналы стандартного вывода}

Корреляция журналов и трейсов реализована в модуле `logging.rs` через форматтер
`JsonTraceFormatter`. Он формирует JSON-строку, в которую записываются время,
уровень, целевой модуль, сведения о сервисе и, при наличии активного контекста,
идентификаторы `trace_id` и `span_id`.

Внутри форматтера поля события сначала собираются во временную структуру, затем
добавляются базовые сведения о сервисе, после чего вызывается процедура
добавления текущего контекста трассировки. За счет этого журнальная запись может
быть сопоставлена с конкретным спаном в системе наблюдаемости.

Представительный вид такой записи можно описать следующим образом:

\begin{verbatim}
{
  "timestamp": "...",
  "level": "INFO",
  "service.name": "checkout-api",
  "deployment.environment": "local",
  "trace_id": "...",
  "span_id": "...",
  "event.name": "checkout.started"
}
\end{verbatim}

Здесь важно отметить, что приведенный пример показывает структуру полей,
формируемых библиотекой, а не конкретную сохраненную строку журнала из
демонстрационного запуска. Фактическую запись можно получить при ручном запуске
примеров или интеграции библиотеки в прикладной сервис.

\section{Пример подключения в Axum}

Практическое подключение библиотеки к Axum подтверждается демонстрационными
сервисами `checkout-api`, `order-processor` и `risk-engine`. Во всех трех
случаях используется одна и та же базовая последовательность:
\begin{enumerate}
    \item построить `TracingConfig`;
    \item вызвать `init_global_tracing`;
    \item создать слой `MyOtelTracingLayer`;
    \item подключить этот слой к `Router`.
\end{enumerate}

Минимальный характерный фрагмент выглядит следующим образом:

\begin{verbatim}
let config = tracing_setup::tracing_config("checkout-api")?;
let _guard = init_global_tracing(config)?;
let layer = tracing_setup::tracing_layer()?;

let app = Router::new()
    .route("/checkout", get(checkout))
    .route("/health", get(health))
    .with_state(state)
    .layer(layer);
\end{verbatim}

Дополнительно `checkout-api` настраивает политику захвата заголовков:
используются стандартные идентификаторы запросов, служебные заголовки шлюза,
идентификатор арендатора и маскируемый идентификатор пользователя. Это
подтверждает, что библиотека рассчитана не только на минимальное подключение,
но и на более строгую эксплуатационную настройку атрибутов запроса.

Для исходящих вызовов сервисы `checkout-api` и `order-processor` используют
`TracedHttpClient`, который оборачивает `reqwest::Client` и автоматически
создает клиентский спан перед вызовом метода `send`. Тем самым библиотека
объединяет серверную и клиентскую стороны трассировки в единый прикладной
контур.

\section{Модульные и интерфейсные тесты}

Автоматические проверки библиотеки расположены в двух группах:
\begin{enumerate}
    \item внутренние тесты внутри исходных файлов библиотеки;
    \item отдельные тесты публичного интерфейса в каталоге
    `crates/my-infra-otel/tests`.
\end{enumerate}

Для проверки фактического состояния проекта была выполнена команда
`cargo test`. Она завершилась успешно и подтвердила:
\begin{enumerate}
    \item 29 внутренних тестов библиотеки;
    \item 18 внешних тестов публичного интерфейса и наблюдаемого поведения;
    \item 1 тест в `demo-ui`;
    \item 5 успешно пройденных примеров из документации.
\end{enumerate}

\begin{table}[H]
    \centering
    \caption{Основные автоматические проверки библиотеки}
    \label{tab:tests-summary}
    \renewcommand{\arraystretch}{1.2}
    \begin{tabularx}{\textwidth}{>{\raggedright\arraybackslash}p{3.0cm} >{\raggedright\arraybackslash}X >{\raggedright\arraybackslash}p{2.5cm} >{\raggedright\arraybackslash}p{1.8cm}}
        \toprule
        Группа тестов & Что проверяется & Основные файлы & Статус \\
        \midrule
        Конфигурация & Значения по умолчанию, типизированные ошибки, валидация адреса OTLP и тайм-аутов & `config.rs`, `tests/config.rs` & Пройдено \\
        \addlinespace
        Захват заголовков & Запрещенные заголовки, дубли правил, режимы записи, ограничения длины и готовые профили захвата & `header_capture.rs`, `tests/header_capture.rs` & Пройдено \\
        \addlinespace
        Атрибуты заголовков & Проверка корректности имени заголовка и ключа атрибута & `labels.rs`, `tests/header_attrs.rs` & Пройдено \\
        \addlinespace
        Передача контекста & Внедрение и извлечение `traceparent`, экспорт клиентского спана, фиксация ошибочного ответа без утечки полного адреса & `propagation.rs`, `tests/propagation.rs` & Пройдено \\
        \addlinespace
        Журналы и события & Добавление `trace_id` и `span_id` в журналы, корректная запись прикладного события без инициализированного подписчика & `logging.rs`, `events.rs`, `tests/log_correlation.rs` & Пройдено \\
        \addlinespace
        Завершение работы & Управляемое завершение и защита от повторной инициализации подсистемы трассировки & `guard.rs`, `tests/shutdown.rs` & Пройдено \\
        \bottomrule
    \end{tabularx}
\end{table}

Таким образом, библиотека уже имеет устойчивое автоматическое покрытие для
основных элементов своего контракта: конфигурации, обработки заголовков,
передачи контекста, корреляции журналов и завершения работы.

\section{Интеграционный сценарий нескольких сервисов}

Помимо модульных и интерфейсных тестов проект содержит воспроизводимый
многосервисный контур в `infra/docker-compose.yml`. В нем поднимаются Jaeger,
OpenTelemetry Collector, а также четыре приложения на Rust:
\begin{enumerate}
    \item `checkout-api`;
    \item `order-processor`;
    \item `risk-engine`;
    \item `demo-ui`.
\end{enumerate}

Этот сценарий проверяет библиотеку в распределенной цепочке, где запрос
попадает в `checkout-api`, затем инициирует прямой вызов `risk-engine`,
вызов `order-processor` и вложенный вызов `risk-engine` из
`order-processor`. Следовательно, один трейс проходит через несколько
независимых HTTP-сервисов и включает как серверные, так и клиентские спаны.

Здесь необходимо явно зафиксировать текущее ограничение проекта. В каталоге
автоматических тестов нет отдельного полноформатного теста, который поднимает
два или более сервиса как независимые процессы и затем автоматически проверяет
единый идентификатор трейса на всем пути запроса. Ближайшим автоматическим
подтверждением служит `tests/propagation.rs`, где на временном HTTP-сервере
проверяется внедрение `traceparent` и корректность клиентского спана. Полный
многосервисный проход в текущем состоянии проекта выступает как
воспроизводимая демонстрационная проверка, а не как отдельный автоматический
интеграционный тест.

\section{Проверка трейса в Jaeger}

Для визуальной проверки результата в проект включена связка OpenTelemetry
Collector и Jaeger. В файле `infra/otel-collector-config.yml` настроен прием
данных по OTLP/HTTP на `0.0.0.0:4318`, пакетная обработка и последующая
передача трейсов в Jaeger. Файл `infra/docker-compose.yml` поднимает контейнеры
`otel-collector` и `jaeger`, а также передает всем демонстрационным сервисам
адрес экспортера через переменную среды `OTLP_ENDPOINT`.

Ожидаемая ручная проверка в Jaeger включает следующие действия:
\begin{enumerate}
    \item запустить инфраструктуру и демонстрационные сервисы;
    \item инициировать запрос к `checkout-api` или сценарий через `demo-ui`;
    \item открыть Jaeger и выбрать последний трейс сервиса `checkout-api`;
    \item убедиться, что в цепочке присутствуют серверный спан входящего
    запроса, клиентские спаны зависимых вызовов и продолжение трейса в
    `order-processor` и `risk-engine`.
\end{enumerate}

По исходным файлам проекта библиотека и демонстрационный контур полностью
подготовлены к такой проверке. Однако в текущем наборе материалов четвертой
главы отсутствует сохраненный снимок экрана Jaeger с конкретным трейсом.
Следовательно, для финальной версии пояснительной записки необходимо добавить
рисунок с изображением фактического трейса и подписью, описывающей цепочку
спанов.

\textbf{FIXME:} после ручного запуска демонстрационного сценария добавить
снимок экрана Jaeger и ссылку на него в тексте главы.

\section{Выводы по реализации и тестированию}

Реализация библиотеки подтверждает основные решения, принятые на этапе
проектирования. Конфигурация сведена к единой модели, промежуточный слой
автоматически создает серверные спаны, клиентская обертка переносит контекст в
исходящих вызовах, а экспорт строится на основе OTLP/HTTP и внешней
инфраструктуры наблюдаемости.

Автоматические тесты подтверждают корректность ключевых механизмов: валидации
конфигурации, безопасного захвата заголовков, передачи контекста, корреляции
журналов и управляемого завершения. Демонстрационный многосервисный контур
показывает применимость библиотеки в распределенном HTTP-сценарии. Для полного
закрытия материалов главы остается добавить визуальное подтверждение трейса в
Jaeger и, при необходимости, в дальнейшем выделить отдельный автоматический
межсервисный тест уровня нескольких процессов.
